\newpage
\phantomsection
\label{prf:one-inverse-is-one}

\begin{remark*}[Return]
\hyperref[thm:one-inverse-is-one]{Return to Theorem}
\end{remark*}

\begin{theorem*}[The Inverse of One Is One]
In any field $F$,
\[
1^{-1} = 1.
\]
\end{theorem*}

\begin{proof}
\textbf{Professional Standard Proof.}~\\
Since $1\cdot 1 = 1$, the element $1$ is a multiplicative inverse of $1$.
By uniqueness of multiplicative inverses, $1^{-1} = 1$.
\end{proof}

\begin{proof}
\textbf{Detailed Learning Proof.}~\\

\textbf{Step 1.} Recall the defining condition for a multiplicative inverse.

An element $y\in F$ is the multiplicative inverse of $x\in F\setminus\{0\}$
if and only if
\[
x\cdot y = 1.
\]
By uniqueness of multiplicative inverses
(\hyperref[thm:uniqueness-of-multiplicative-inverse]{Theorem~\ref*{thm:uniqueness-of-multiplicative-inverse}}),
there is exactly one such $y$, denoted $x^{-1}$.

\textbf{Step 2.} Verify that $1$ satisfies the defining condition for $1^{-1}$.

Compute:
\[
1\cdot 1 = 1.
\]
This uses the multiplicative identity axiom: $1$ is the identity element,
so $1\cdot 1 = 1$.

\textbf{Step 3.} Conclude by uniqueness.

The equation $1\cdot 1 = 1$ states that $1$ is a multiplicative inverse
of $1$.  By uniqueness of multiplicative inverses,
\[
1^{-1} = 1. \qedhere
\]
\end{proof}

\begin{remark*}[Proof structure]
The proof is a one-line uniqueness argument.  Exhibit $1$ as a solution to
$1\cdot y = 1$, which holds trivially by the multiplicative identity axiom.
Uniqueness of multiplicative inverses then pins down $1^{-1} = 1$.  No field
axiom beyond the multiplicative identity and uniqueness of inverses is
required.
\end{remark*}

\begin{remark*}[Dependencies]~\\
\begin{itemize}
  \item \hyperref[def:field]{Definition~\ref*{def:field}}
        \textnormal{(multiplicative identity axiom: $1\cdot 1 = 1$)}
  \item \hyperref[thm:uniqueness-of-multiplicative-inverse]{Theorem~\ref*{thm:uniqueness-of-multiplicative-inverse}}
        \textnormal{(the multiplicative inverse of any nonzero element is unique)}
\end{itemize}
\end{remark*}
